\DocumentMetadata{
  pdfversion=2.0,
  pdfstandard=ua-2,
  lang=en-US,
  tagging=on,
  tagging-setup={math/setup=mathml-SE,tabsorder=structure}
}
\documentclass[11pt,letterpaper]{article}
\usepackage[margin=0.68in,includefoot]{geometry}
\usepackage{newcomputermodern}
\usepackage{amsmath,amscd,mathtools}
\usepackage{tikz,adjustbox}
\providecommand{\square}{\mdlgwhtsquare}
\usepackage[hidelinks]{hyperref}
\hypersetup{
  pdftitle={Algebra First-Year Exam, August 2019, Part 1},
  pdfauthor={University of Florida Department of Mathematics},
  pdfsubject={Graduate mathematics examination},
  pdfdisplaydoctitle=true,
  bookmarksopen=true
}
\setcounter{secnumdepth}{0}
\setlength{\parindent}{0pt}
\setlength{\parskip}{0.28em}
\setlength{\emergencystretch}{3em}
\setlength{\leftmargini}{2.15em}
\setlength{\leftmarginii}{2.65em}
\setlength{\leftmarginiii}{3em}
\renewcommand{\labelenumi}{\textbf{\theenumi.}}
\pagestyle{empty}
\raggedbottom
\allowdisplaybreaks
\newenvironment{examproblems}[1]{%
  \begin{enumerate}%
  \setcounter{enumi}{#1}%
  \setlength{\topsep}{0.35em}%
  \setlength{\partopsep}{0pt}%
  \setlength{\itemsep}{0.48em}%
  \setlength{\parsep}{0.18em}%
}{\end{enumerate}}
\newenvironment{examsubparts}{%
  \begin{enumerate}%
  \setlength{\topsep}{0.18em}%
  \setlength{\partopsep}{0pt}%
  \setlength{\itemsep}{0.16em}%
  \setlength{\parsep}{0.1em}%
}{\end{enumerate}}
\newenvironment{examinstructions}{%
  \par\begingroup\itshape
}{%
  \par\endgroup\vspace{0.18em}
}
\makeatletter
\renewcommand\section{\@startsection{section}{1}{0pt}%
  {-1.25ex plus -.25ex minus -.1ex}%
  {0.55ex plus .1ex}%
  {\normalfont\large\bfseries}}
\renewcommand{\@maketitle}{%
  \newpage\null\vskip -1em%
  {\footnotesize\bfseries
    \href{https://www.ufl.edu/}{University of Florida}, \href{https://math.ufl.edu/}{Department of Mathematics}\par}%
  \vskip -0.5em%
  \tagstructbegin{tag=Title}%
    {\LARGE\bfseries\href{https://gma.math.ufl.edu/exams/algebra-first-year/}{Algebra First-Year Exam}, August 2019, Part 1\par}%
  \tagstructend%
  \vskip 0.25em%
  \tagmcbegin{artifact}%
    \hrule height 1pt%
  \tagmcend%
  \vskip 1em%
}
\makeatother
\title{Algebra First-Year Exam, August 2019, Part 1}
\author{}
\date{}
\begin{document}
\maketitle
\thispagestyle{empty}
\begin{examinstructions}
Answer four problems. (If you turn in more, the first four will be graded.) Within reason, you may quote theorems as long as you state them clearly.
\end{examinstructions}
\begin{examproblems}{0}
\item
Let~\(p\) be a prime, let~\(G\) be a finite group, and let~\(P\) be a Sylow~\(p\)-subgroup of~\(G\). Let~\(H\) be a subgroup of~\(G\) containing the normalizer \(N_G(P)\) of~\(P\) in~\(G\). Prove that then \(N_G(H)=H\).
\item
State and prove the Second Isomorphism Theorem for groups.
\item
Let \(n\geq 2\), and let \(S_n\) be the symmetric group on~\(n\) letters.
\begin{examsubparts}
\item[\textnormal{(a)}]
Define what is a transposition of \(S_n\).
\item[\textnormal{(b)}]
Give an example of a transposition of \(S_n\).
\item[\textnormal{(c)}]
Prove that every element of \(S_n\) is a product of transpositions of \(S_n\).
\end{examsubparts}
\item
Let~\(G\) be a finite group, and assume that~\(G\) acts transitively on the finite non-empty set~\(\Omega\). Let \(\omega\in\Omega\). We denote by \(G_\omega\) the set of elements of~\(G\) that fix~\(\omega\).
\begin{examsubparts}
\item[\textnormal{(a)}]
Prove that \(G_\omega\) is a subgroup of~\(G\).
\item[\textnormal{(b)}]
Prove that the number of elements in~\(\Omega\) is exactly \([G:G_\omega]\).
\end{examsubparts}
\item
Let~\(G\) be a finite group.
\begin{examsubparts}
\item[\textnormal{(a)}]
Define what it means to say that~\(G\) is solvable.
\item[\textnormal{(b)}]
Without assuming any result about solvable groups (i.e. just from your definition), prove that if the order of~\(G\) is \(20\) then~\(G\) is solvable.
\end{examsubparts}
\item
Prove or disprove the following statement. Let~\(n\) be a natural number, let \(S_n\) be the symmetric group on~\(n\) letters, let \(A_n\) be the alternating group on~\(n\) letters, and let \(\sigma\in S_n\). Then, if~\(\sigma\) has order \(2019\) then \(\sigma\in A_n\).
\end{examproblems}
\end{document}
